\section{Related Works}
\subsection{Document-level Relation Extraction.}
\label{subsec:docre-models-related-works}
Document-level Relation Extraction (DocRE) identifies relations between entity pairs across sentences or entire documents. Existing methods mainly fall into two categories: \textit{transformer-based} and \textit{graph-based}. Transformer-based approaches \citep{zhou2021document, tan2022document, zhang2021document} enhance entity representations using attention mechanisms in pre-trained language models, including localised context embeddings \citep{zhang2021document} and multi-hop reasoning via Axial Attention \citep{tan2022document}. In contrast, graph-based methods \citep{christopoulou2019connecting, zeng2020double, xu2021document} model document structure explicitly by constructing heterogeneous graphs (e.g., mention, entity, and sentence nodes) and applying graph-based reasoning to improve relation prediction.

\subsection{Synthetic Data for DocRE.}
% Aforementioned DocRE systems in Section \ref{subsec:docre-models-related-works} require substantial amount of high-quality labelled data to achieve satisfactory performance. However, high-quality labelled data is high cost and error-prone due to the nature of manual annotation. Thus, recent work has explored synthetic data to improve scalability and performance. One line of research focuses on generating additional training data \citep{tan2022document, sun2023uncertainty, feng2024dure, tran2025vaediff}. For example, \citet{feng2024dure} propose a text generation approach to create pseudo-texts for long-tail relations, while \citet{tran2025vaediff} introduce a representation-level augmentation method to synthesise pseudo entity pair embeddings. Another direction centres on pseudo-labelling. For instance, \citet{tan2022document} present a knowledge distillation framework in which a teacher model labels unlabelled data to train a student model. Building on this, \citet{sun2023uncertainty} incorporate Monte Carlo dropout and a multi-stage training scheme to assess pseudo-label quality via uncertainty estimation.

% However, most of the proposed methods are computationally expensive because of massive number of documents and entity pairs. \citet{feng2024dure}'s method tried to generate each pseudo-document for each entity pair. While proposed solution \cite{sun2023uncertainty} requires multiple training and denoising rounds to generate reliable pseudo-labels. Therefore, the proposed \model\ method aims to reduce the computational cost with more reliable pseudo-labels compared to previous methods.
Aforementioned DocRE systems in Section \ref{subsec:docre-models-related-works} typically require large amounts of high-quality labelled data, which is costly and prone to annotation errors. To mitigate this, recent work explores synthetic data to improve scalability. One direction focuses on data generation \citep{tan2022document, sun2023uncertainty, feng2024dure, tran2025vaediff}, such as generating pseudo-texts for long-tail relations \citep{feng2024dure} or synthesising entity pair representations \citep{tran2025vaediff}. Another direction leverages pseudo-labelling, where teacher models annotate unlabelled data for student training \citep{tan2022document}, with extensions incorporating uncertainty estimation via Monte Carlo dropout and multi-stage training \citep{sun2023uncertainty}.

However, these methods still have notable limitations. \citet{tran2025vaediff}'s method continues to rely on substantial amounts of high-quality labelled data for strong performance. Meanwhile, methods proposed by \citet{sun2023uncertainty, feng2024dure} are computationally expensive due to the large number of documents and entity pairs. For instance, \citet{feng2024dure} generate pseudo-documents for each entity pair, while \citet{sun2023uncertainty} require multiple rounds of training and denoising. Consequently, \model\ is proposed to reduce computational cost while producing more reliable pseudo-labels than prior methods.